% GENERATED FILE. Edit canonical lessons, then run npm run generate:books.
% canonical-source-hash: fnv1a64:facc0883c4ec7149
% canonical-lessons: ZH-C01-tones, ZH-C01-ni, ZH-C01-hao, ZH-C01-nihao, ZH-C01-tone-sandhi, ZH-C01-hao-fond, ZH-C01-practice

\chapter{Nǐ Hǎo — Hello, Character by Character}
\label{ch:zh-nihao}
\hlchaptermodality{1}

\begin{chapteropening}
\textbf{By the end of this chapter:} \emph{I can greet someone in Mandarin and answer when greeted, reading both characters unaided and saying the greeting with the tone a speaker actually uses rather than the one the dictionary prints.}

A two-line meeting in which both voices say the same word — read from the characters alone, with the first third tone raised by sandhi.
\end{chapteropening}

\section[mā má mǎ mà]{mā má mǎ mà — one syllable, four words}

\label{lesson:ZH-C01-tones}



\begin{quote}
Nothing to recall yet. Before the first Chinese word arrives, one thing your voice is about to be asked to do that English never asks of it.
\end{quote}

\begin{sounds}
In English, pitch carries mood. \emph{Really.} and \emph{Really?} are the same word, said two ways.

In Mandarin, \textbf{pitch is part of the word}, exactly the way a consonant is. Take the syllable \emph{ma} and give it four different pitches and you get four unrelated words:

\begin{itemize}
\raggedright
  \item \textbf{mā} — held high and flat, one steady note. \emph{Mother.}
  \item \textbf{má} — rising, like the pitch on a surprised English \textquotedblleft{}huh?\textquotedblright{}. \emph{Hemp.}
  \item \textbf{mǎ} — dipping low, a little creaky. \emph{Horse.}
  \item \textbf{mà} — falling sharply, like a clipped \textquotedblleft{}No!\textquotedblright{}. \emph{To scold.}
\end{itemize}

There is a fifth, the \textbf{neutral} tone: short, light, no mark written at all.

The accent over the vowel is how the romanization writes the tone: a flat bar in \textbf{ā}, a rising stroke in \textbf{á}, a hook in \textbf{ǎ}, a falling stroke in \textbf{à}. The Chinese characters themselves record no tone anywhere, so the accent is a teaching aid, not part of the writing.

Say all four in order now, slowly: \textbf{mā, má, mǎ, mà}. {[}REPEAT x2{]}
\end{sounds}

\subsection*{Writing — trace four visible paths}
Keep the model visible. With a finger first, then a pencil, trace only the mark above each \textbf{a}: \textbf{ā, á, ǎ, à}. Move flat, up, down-then-up, and down. Say the matching pitch while your pencil moves.

Stop after one slow pass. This is not a memory test, and these pinyin marks are not Chinese characters. The next lesson will give your hand its first two-stroke character; this tiny step only lets pitch become a line you can see and follow.

\subsection*{Your turn}
\begin{itemize}
\raggedright
  \item \emph{Say it:} \textbf{mā} high and flat, then \textbf{mà} falling — hear how far apart they are
  \item \emph{Answer:} if you change the pitch of a Mandarin syllable, what changes?
\end{itemize}

\begin{tcolorbox}[breakable,colback=teal!4,colframe=teal!35!black,title={Before you move on}]
Which accent tells you to dip low? (The hook, as in \textbf{ǎ}.) Does changing a Mandarin tone change how you feel about a word, or which word it is? (Which word it is.)

Source: the tone inventory this track keeps
\end{tcolorbox}
\section[nǐ]{\zh{你} — \textquotedblleft{}you\textquotedblright{}, and the first character}

\label{lesson:ZH-C01-ni}



\begin{quote}
Say the four tones on \emph{ma} once, in order. Hold the third one — low and dipping. The first Chinese word carries exactly that pitch.
\end{quote}

\begin{sounds}
\textbf{nǐ}. The \emph{n} is the English \emph{n}; the \emph{i} is the \emph{ee} of \textquotedblleft{}see\textquotedblright{}. The hook over it is the \textbf{third tone}: start low, stay low, let it creak. {[}REPEAT x2{]}
\end{sounds}

\subsection*{Script — \zh{你}, and the person hiding on its left}
Chinese is not written with letters. It is written with \textbf{characters}, each an evenly sized square block, and each block built from smaller recurring pieces called \textbf{components}.

\textbf{\zh{你}} has two of them, left and right:

\begin{quote}
\textbf{\zh{亻}} + \textbf{\zh{尔}} $\to$ \textbf{\zh{你}}
\end{quote}

\textbf{\zh{亻}} is the \textquotedblleft{}person\textquotedblright{} component. On its own the character for a person is \textbf{\zh{人}} — two legs walking. Push it to the left edge of a block and it is squashed upright into \textbf{\zh{亻}}: one short slant over one vertical. Same piece, standing sideways to make room.

\textbf{\zh{尔}} on the right is there for its \emph{sound}, not its meaning. Chinese builds most of its characters this way — one component pointing at the meaning, another at the pronunciation.

Stroke order is a real, taught system here: top before bottom, left before right. So \textbf{\zh{亻}} first, then \textbf{\zh{尔}}.

\begin{cousinweb}
\textbf{\zh{你}} \emph{nǐ} — \textbf{you}, addressing one person.

Here this book has to be honest with you. A related language lets a new word be anchored to English words you already own through a shared ancestor — Spanish \emph{gracias} to \emph{grace} and \emph{gratitude}, all from Latin \emph{gratia}. \textbf{Chinese shares no ancestor with English.} There is no cousin web for \emph{nǐ}, and inventing one would be worse than having none.

What replaces it is what you just read: the character is itself decomposable, and its pieces mean things. \textbf{\zh{你}} carries a person on its left because it is about a person. That is the hook — not history you share with English, but structure you can see.
\end{cousinweb}

\subsection*{Your turn}
\begin{itemize}
\raggedright
  \item \emph{Say it:} \textbf{nǐ} — low and dipping, not flat
  \item \emph{Read it:} \textbf{\zh{你}} — name its left piece, then its right piece
\end{itemize}

\begin{tcolorbox}[breakable,colback=teal!4,colframe=teal!35!black,title={Before you move on}]
Which component of \textbf{\zh{你}} tells you the word is about a person? (\textbf{\zh{亻}}.)

Source: \href{https://www.unicode.org/charts/PDF/U4E00.pdf}{Unicode CJK Unified Ideographs chart}
\end{tcolorbox}
\section[hǎo]{\zh{好} — \textquotedblleft{}good\textquotedblright{}, and what replaces etymology here}

\label{lesson:ZH-C01-hao}



\begin{quote}
Say the word for \textquotedblleft{}you\textquotedblright{} without looking. Which of its two components is the person? Now hold that low, dipping pitch — the next word wants the same one.
\end{quote}

\begin{sounds}
\textbf{hǎo}. The \emph{h} is a little rougher than English \emph{h}, scraped at the back of the mouth. The \emph{ao} is one gliding vowel, the \emph{ow} of \textquotedblleft{}how\textquotedblright{}. The hook accent is the same third tone you already met: low, dipping, creaky. {[}REPEAT x2{]}
\end{sounds}

\subsection*{Script — \zh{好}, a woman beside a child}
Two components again, side by side:

\begin{quote}
\textbf{\zh{女}} + \textbf{\zh{子}} $\to$ \textbf{\zh{好}}
\end{quote}

\textbf{\zh{女}} on the left is the character for \textquotedblleft{}woman\textquotedblright{}: a crossing stroke, a curved leg, a horizontal through the middle. \textbf{\zh{子}} on the right is \textquotedblleft{}child\textquotedblright{}: a hooked top, a hooked vertical, a horizontal arm. Left before right, as always.

Notice what just happened. \textbf{\zh{女}} and \textbf{\zh{子}} are not letters spelling out the sound \emph{hǎo} — neither is pronounced anything like it. They are whole characters with their own meanings, reused as building blocks.

\begin{cousinweb}
\textbf{\zh{好}} \emph{hǎo} — \textbf{good}, and also \textbf{well}.

A woman beside a child, and the character means \textquotedblleft{}good\textquotedblright{}. That picture is the hook — and it stands in for the cousin web that a related language would supply, because \textbf{Chinese and English share no ancestor} and there is no honest chain from \emph{hǎo} to any English word.

Two cautions, because a memory aid that pretends to be history rots:

\begin{itemize}
\raggedright
  \item \textquotedblleft{}Woman + child = good\textquotedblright{} is the \textbf{traditional gloss}. It is genuinely useful for remembering the shape, but palaeographers do not all accept it as the historical origin; some read \textbf{\zh{子}} as a phonetic element instead. A picture that helps, not a proven derivation.
  \item Component glosses are hooks, not translations. \textbf{\zh{好}} does not \textquotedblleft{}mean\textquotedblright{} woman-child. It means good.
\end{itemize}
\end{cousinweb}

\subsection*{Your turn}
\begin{itemize}
\raggedright
  \item \emph{Say it:} \textbf{hǎo} — low and dipping, the \emph{ow} of \textquotedblleft{}how\textquotedblright{}
  \item \emph{Read it:} \textbf{\zh{好}} — name its left component, then its right one
  \item \emph{Answer:} why does this track use component pictures instead of English cousins?
\end{itemize}

\begin{tcolorbox}[breakable,colback=teal!4,colframe=teal!35!black,title={Before you move on}]
What does \textbf{\zh{好}} mean? (Good, or well.) Is its woman-and-child picture a proven etymology? (No — a traditional and useful gloss, disputed as history.)

Source: \href{https://www.unicode.org/charts/PDF/U4E00.pdf}{Unicode Han database, character \zh{好}}
\end{tcolorbox}
\section[nǐ hǎo]{\zh{你好} — the greeting, assembled from two things you already have}

\label{lesson:ZH-C01-nihao}



\begin{quote}
Two characters, both already yours. Say the one that means \textquotedblleft{}you\textquotedblright{}. Say the one that means \textquotedblleft{}good\textquotedblright{}. Now put them in that order and you have said hello — but say it slowly for one more moment, because the tones are about to do something.
\end{quote}

\subsection*{The exchange}
\begin{quote}
\textbf{\zh{你好}} \emph{nǐ hǎo!} Hello.
\end{quote}

\begin{quote}
\textbf{\zh{你好}} \emph{nǐ hǎo!} Hello.
\end{quote}

The greeting and its answer are the same word. Say it to one person you have just met, at any hour of the day.

\begin{grammarlens}[title={Grammar Lens — a word can be more than one character}]
English counts a word in \textbf{letters}. Chinese does not have letters, so the tempting guess is that it counts a word in \textbf{characters} — one character, one word. That guess is wrong, and correcting it early saves a great deal of confusion.

A character writes one \textbf{syllable} and usually one \textbf{meaningful piece}. A \textbf{word} is one or more of those pieces:

\begin{itemize}
\raggedright
  \item \textbf{\zh{你}} — one character, one word: \emph{you}
  \item \textbf{\zh{好}} — one character, one word: \emph{good}
  \item \textbf{\zh{你好}} — two characters, still \textbf{one word}: \emph{hello}
\end{itemize}

Nobody hearing \textbf{\zh{你好}} parses it as \textquotedblleft{}you good\textquotedblright{}, any more than an English speaker hears \textquotedblleft{}good\textquotedblright{} and \textquotedblleft{}bye\textquotedblright{} inside \emph{goodbye}. The pieces are visible, which makes the word easy to remember; the whole is what it means.
\end{grammarlens}

\begin{culture}
\textbf{\zh{你好}} is safe, neutral, and understood everywhere — which is exactly why a course teaches it first. But it is worth knowing that it is used less freely than English \textquotedblleft{}hello\textquotedblright{}. Between people who already know each other, a Mandarin speaker is more likely to open with a small situational remark than with \textbf{\zh{你好}} at all. Reserve it for meeting someone, and you will sound right.
\end{culture}

\subsection*{Your turn}
\begin{itemize}
\raggedright
  \item \emph{Say it:} \textbf{nǐ hǎo} — greeting someone you have just met
  \item \emph{Say it:} \textbf{nǐ hǎo} — answering that same greeting
  \item \emph{Answer:} how many characters is this word, and how many words is it?
\end{itemize}

\begin{tcolorbox}[breakable,colback=teal!4,colframe=teal!35!black,title={Before you move on}]
Which character comes first, and what does it mean on its own? (\textbf{\zh{你}}, \textquotedblleft{}you\textquotedblright{}.) Is \textbf{\zh{你好}} one word or two? (One.)

Source: \href{https://www.unicode.org/charts/PDF/U4E00.pdf}{Unicode Han database, characters \zh{你} and \zh{好}}
\end{tcolorbox}
\section[ní hǎo]{ní hǎo — why the greeting is not said the way it is written}

\label{lesson:ZH-C01-tone-sandhi}



\begin{quote}
Say the greeting once, slowly, giving both syllables the low dipping third tone you were taught. It should feel awkward in the mouth. That awkwardness is the whole lesson.
\end{quote}

\begin{grammarlens}[title={Grammar Lens — one rule, and it is not optional}]
\textbf{When two third tones come one after the other, the first is spoken as a second (rising) tone.}

So the greeting, written \textbf{nǐ hǎo}, is actually said \textbf{ní hǎo} — the first syllable climbs instead of dipping. {[}REPEAT x2{]}

Nothing about this shows in the writing. Dictionaries print \textbf{nǐ hǎo}, because that is the citation tone each word carries on its own, and the characters \textbf{\zh{你好}} never change at all. A learner who trusts only what is written will say the commonest greeting in the language wrong every time.

Why does it happen? Two dipping tones in a row would need the voice to sink, climb, sink and climb again inside a fifth of a second. Speakers do not do that; the first one gives way. This is a real rule of the sound system, not a sloppy shortcut, and every native speaker applies it without thinking.

Compare it to what English does with \emph{\textquotedblleft{}the\textquotedblright{}}: written one way, said \emph{thuh} before a consonant and \emph{thee} before a vowel, and no English speaker thinks of that as two words. The difference is that here the change lands on \textbf{pitch} — the same property that distinguishes one word from another.
\end{grammarlens}

\subsection*{Your turn}
\begin{itemize}
\raggedright
  \item \emph{Say it:} \textbf{ní hǎo} — first syllable rising, second dipping
  \item \emph{Say it:} it again, at speaking speed, without the pause between them
  \item \emph{Answer:} which spelling changes when the rule applies — pinyin, characters, or neither?
\end{itemize}

\begin{tcolorbox}[breakable,colback=teal!4,colframe=teal!35!black,title={Before you move on}]
Two third tones meet. What happens to the first? (It rises.) What happens to the characters? (Nothing.)

Source: \href{https://ctext.org/dictionary.pl}{Chinese Text Project, on tone in Mandarin}
\end{tcolorbox}
\section[hào]{\zh{好} again — proof, on a character you own, that tone is the word}

\label{lesson:ZH-C01-hao-fond}



\begin{quote}
Name the character built from a woman and a child, and say what it means. Hold its tone: low and dipping.
\end{quote}

\begin{cousinweb}
The very same character, \textbf{\zh{好}}, has a second reading — and the only thing that separates the two is the tone:

\begin{itemize}
\raggedright
  \item \textbf{\zh{好}} \emph{hǎo}, dipping third tone — \textbf{good}
  \item \textbf{\zh{好}} \emph{hào}, sharply falling fourth tone — \textbf{to be fond of, to like}
\end{itemize}

Same strokes. Same components. Same shape on the page. Two words.

This is the cleanest proof available of what the first lesson claimed. Tone is not decoration on a word; it is \emph{inside} the word, doing the job a consonant does in English. Change \emph{hǎo} to \emph{hào} and you have not said \textquotedblleft{}good\textquotedblright{} with feeling — you have said a different word, the way \emph{bat} and \emph{pat} are different words.

The two readings are related in meaning, and the relation is a good hook: what is \textbf{good} is what one is \textbf{fond of}. That is a link inside Chinese, between two readings of one character. It is not a link to English. There is none.

You are not asked to produce \textbf{hào} yet. Recognising it is the point: when a familiar character turns up wearing an unfamiliar tone, the tone is telling you this is a different word.
\end{cousinweb}

\subsection*{Your turn}
\begin{itemize}
\raggedright
  \item \emph{Answer:} \textbf{hǎo} or \textbf{hào} — which one means \textquotedblleft{}good\textquotedblright{}?
  \item \emph{Answer:} what changed between them, and what stayed the same?
  \item \emph{Answer:} which two components is this character built from?
\end{itemize}

\begin{tcolorbox}[breakable,colback=teal!4,colframe=teal!35!black,title={Before you move on}]
One character carried two words today. What told them apart? (Only the tone.)

Source: \href{https://www.unicode.org/charts/PDF/U4E00.pdf}{Unicode Han database, character \zh{好}}
\end{tcolorbox}
\section[Practice]{Practice — greet someone, and be greeted back}

\label{lesson:ZH-C01-practice}



\begin{quote}
Retrieve the two words separately first: the one meaning \textquotedblleft{}you\textquotedblright{}, then the one meaning \textquotedblleft{}good\textquotedblright{}. Do not run them together yet.
\end{quote}

\subsection*{The exchange}
\begin{quote}
A: \textbf{\zh{你好}} \emph{ní hǎo!} B: \textbf{\zh{你好}} \emph{ní hǎo!}
\end{quote}

Two lines, one word, both voices. Every piece is already yours: the characters from their own lessons, the meaning from the greeting lesson, the rising first syllable from the sandhi rule. Nothing new is hiding here.

Say A's line, pause, then answer as B. {[}REPEAT x2{]}

\begin{grammarlens}[title={What you've built this chapter}]
Four things, and only one of them is vocabulary:

\begin{itemize}
\raggedright
  \item \textbf{\zh{你}} is \zh{亻} plus \zh{尔}; \textbf{\zh{好}} is \zh{女} beside \zh{子}. Characters are built from parts.
  \item Those parts are the memory hook this track uses, because Chinese and English share no ancestor and therefore no cousin web.
  \item \textbf{\zh{你好}} is two characters and one word. Characters are not letters.
  \item Pitch is part of the word, and a neighbouring word can change it without changing a stroke.
\end{itemize}
\end{grammarlens}

\subsection*{Your turn}
\begin{itemize}
\raggedright
  \item \emph{Say it:} greet someone you have just met
  \item \emph{Say it:} answer that greeting
  \item \emph{Read it:} \textbf{\zh{你好}} aloud from the characters alone, with no pinyin
\end{itemize}

\begin{tcolorbox}[breakable,colback=teal!4,colframe=teal!35!black,title={Before you move on}]
Run both lines once more. If the first syllable dips instead of rising, go back to the sandhi lesson only — not to the whole chapter.

Source: \href{https://www.unicode.org/charts/PDF/U4E00.pdf}{Unicode Han database, characters \zh{你} and \zh{好}}
\end{tcolorbox}
